\section{Розділ~\ref{sec:livingmachine}. Машина з живих нейронів}

Поведінку культур на ґратках електродів виміряно в прямих дослідах із записом і стимуляцією; механізм залпів спирається на модель виснаження синапсів і досліди на мікрокультурах, а твердження про дофамін у чашці --- на поодинокі досліди, частину яких самі автори вважають лише демонстрацією.

{\footnotesize
\setlength{\tabcolsep}{3pt}\renewcommand{\arraystretch}{1.15}
\begin{longtable}{>{\raggedright}p{0.5cm}>{\raggedright}p{4.0cm}>{\raggedright}p{5.6cm}>{\raggedright}p{2.3cm}>{\raggedright\arraybackslash}p{1.7cm}}
\toprule
№ & Твердження & Дослід і що виміряно & Джерело & Статус \\
\midrule\endfirsthead
\toprule
№ & Твердження & Дослід і що виміряно & Джерело & Статус \\
\midrule\endhead
1 & Культура на чипі: здебільшого \nt{GLUT}, менша частка \nt{GABA}, що залежить від препарату; у культурах гіпокампа близько $6\%$ & Мережа з тисяч нейронів на ґратці електродів --- стандартний дослід (рядок~3). Частку \nt{GABA}-нейронів у культурах гіпокампа рахували за фарбуванням на GABA: близько $6\%$ нейронів. Для культур кори перевіреного числа ми не наводимо & \cite{Benson1994} & виміряно \\
2 & Органоїди: тривимірна нервова тканина розміром близько міліметра зі стовбурових клітин людини & Зі стовбурових клітин виростили тривимірні органоїди з кількома ділянками мозку, зокрема корою з шарами попередників і зрілими нейронами & \cite{Lancaster2013} & виміряно \\
3 & Без входу культура спалахує залпами з інтервалами від секунди до хвилин, залежно від культури & $58$ культур кори щільністю від $3\,000$ до $50\,000$ нейронів записували п'ять тижнів ($963$ півгодинні записи): після двох тижнів майже в усіх культурах переважають залпи всієї мережі; інтервали між залпами --- від $1$ до $300$~с і сильно залежать від препарату & \cite{Wagenaar2006} & виміряно \\
4 & Залп закінчується виснаженням медіатора, інтервал $T_{\text{залп}}\approx\tau_{\text{відн}}\ln\frac{1}{1-x^*}$~\eqref{eq:burstinterval}; $2$--$4$~с --- оцінка моделі з $\tau_{\text{відн}}=0.8$~с, виміряне відновлення повільніше & Формулу виведено в розділі. Модель: мережа з синапсами, що виснажуються, сама породжує залпи всієї мережі. Дослід на мікрокультурах із $4$--$30$ нейронів: тривалість залпу залежить від часу після попереднього, виснаження пухирців медіатора скасовує залпи; висновок авторів --- залп обмежений запасом медіатора. Відновлення запасу там --- десятки секунд, що з тією самою формулою дає інтервали в десятки секунд і хвилини & висновок у розділі; \cite{Tsodyks2000,CohenSegal2011,Tsodyks1997} & теорія \\
5 & «Учитель, який замовкає»: стимуляцію вимикають, коли з'явилася потрібна відповідь, і відповідь закріплюється & Культуру стимулювали з частотою $0.3$--$1$~Гц, доки через $50\pm10$~мс після стимулу не з'явиться обрана відповідь, потім стимуляцію вимикали на $5$~хв; після кількох циклів потрібна відповідь викликалася одразу; будували криві навчання & \cite{Shahaf2001} & виміряно \\
6 & Культура грає в теніс; значуще зростання серій за сеанс --- тільки за повного зворотного зв'язку & Докладно --- у розділі~\ref{sec:dishbrain} і додатку до нього: коли замість стимулу тиша, культура наприкінці сеансу теж краща, ніж без зворотного зв'язку, але з меншим ефектом; навчання від сеансу до сеансу нестійке & \cite{Kagan2022} & виміряно \\
7 & Те, що культура вчиться передбачати свій вхід, --- гіпотеза; гучні слова про «розумність» оскаржено & Принцип вільної енергії --- теоретична пропозиція; прямого досліду, що відрізняв би його від простіших пояснень, немає. Тридцять дослідників у статті-відповіді вказали, що зміна поведінки в петлі не показує ні розуму, ні відчуттів & \cite{Friston2010,Balci2023} & гіпотеза \\
8 & Культура веде віртуальне тіло до цілі за дібраного малюнка стимулів & Програма сама добирала навчальні малюнки стимуляції за поточним результатом; за десятки хвилин мережа досягала заданих станів, пластичність трималася $80$~хв після $2$~год навчання. Ті самі стимули, подані без зв'язку з поведінкою, не діяли & \cite{Bakkum2008} & виміряно \\
9 & Резервуар: навчається тільки лінійне зчитування $y=W_{\text{вих}}\mathbf r$~\eqref{eq:reservoir} & Теорія обчислення на резервуарі: будь-яка нелінійна система зі згасною пам'яттю плюс навчений лінійний вихід & \cite{Maass2002,JaegerHaas2004} & теорія \\
10 & Органоїд як резервуар розрізняє мовців із точністю близько $78\%$ (за даними авторів); за помилкою навчається зчитування, а зв'язки органоїду змінюються без учителя & Звуки мовлення кількох мовців подавали органоїду малюнками струму на ґратці електродів; навчене зчитування розрізняло мовця з точністю близько $78\%$ --- так повідомляють автори. Вони ж повідомляють, що під час тренування змінювалася функціональна зв'язність органоїду, і називають це навчанням без учителя в самому органоїді & \cite{Cai2023} & виміряно \\
11 & Мільйон нейронів витрачає на імпульси близько $10^{-4}$~Вт~\eqref{eq:dishpower} & Оцінка: ціна імпульсу $10^{-10}$~Дж з енергетичного бюджету кори~\eqref{eq:energy}, помножена на кількість імпульсів; прямого вимірювання потужності культури немає & висновок у розділі; \cite{Attwell2001} & теорія \\
12 & Дофамін за спалахом світла: $1$~мМ дофаміну в клітці, $240$ повторів, кількість викликаних імпульсів зростала & На віддаленій платформі з органоїдами дофамін у світлочутливій клітці ($1$~мМ) вивільняли ультрафіолетом ($365$~нм, $0.8$~с), коли стимул викликав більше імпульсів; за $240$ кроків (близько $5$~год) кількість імпульсів загалом зростала. Автори пишуть, що за одним дослідом зв'язок із дофаміном встановити не можна; контролів немає & \cite{Jordan2024} & гіпотеза \\
13 & Дофамін від живих клітин змінює вагу органічного транзистора надовго й оборотно & Клітини, що виділяють дофамін, вирощено над органічним нейроморфним елементом; окиснення дофаміну біля електрода змінює провідність каналу; показано довготривалу зміну ваги та її повернення & \cite{Keene2020} & виміряно \\
14 & Органоїди з робочими \nt{DOPA}-нейронами існують; як учителя їхній дофамін не використовували & В органоїдах зі стовбурових клітин людини знайдено електрично активні зрілі \nt{DOPA}-нейрони й вироблення дофаміну. Дослідів, де цей дофамін навчає іншу мережу, ми в літературі не знайшли & \cite{Jo2016} & виміряно \\
15 & Органоїд балансує жердиною: навчальні стимули від програми кращі за випадкові, без закріплення, через \nt{GLUT}-синапси & Органоїди кори миші в замкненій петлі із задачею «жердина на візку»: у більшості органоїдів сигнали, обрані навчанням із підкріпленням, дали кращий результат, ніж випадкові або жодні; після $45$~хв відпочинку покращення не зберігалося; блокада AMPA- і NMDA-рецепторів його скасовувала & \cite{Robbins2026} & виміряно \\
16 & Без регістрів керування мережа на стенді не може сама вирішити, що вважати успіхом (твердження~\ref{prop:livingmachine}) & В усіх дослідах рядків 5--15 навчальний сигнал задає експериментатор або програма; досліду, де культура сама виробила б критерій успіху, немає --- це висновок із відсутності, а не дослід & --- & гіпотеза \\
\bottomrule
\end{longtable}}
